\section{Hệ thống chiều khái niệm}\label{sec:chieu}
\subsubsection*{Bảng chiều (dimension)}
Kho dữ liệu có 10 bảng chiều. Mỗi bảng có khoá thay thế dạng số nguyên tự tăng, khoá tự nhiên ràng buộc duy nhất, và một dòng Unknown với khoá -1 để bảng fact không bao giờ có khoá rỗng. Chiều thời gian \code{dim\_\allowbreak{}date} phủ 1.461 ngày từ 2016 đến 2019 với ngày lễ Brazil và Black Friday, và được dùng lại cho nhiều mốc thời gian của đơn hàng qua 11 view đóng vai. Chiều địa lý \code{dim\_\allowbreak{}geography} có một dòng cho mỗi mã bưu chính, hợp nhất từ bảng toạ độ, khách hàng và người bán, với phân cấp Zip $\rightarrow$ Thành phố $\rightarrow$ Bang $\rightarrow$ Vùng. Chiều khách hàng \code{dim\_\allowbreak{}customer} áp dụng SCD loại 2: khi một khách đổi mã bưu chính qua các đơn, hệ thống đóng phiên bản cũ và mở phiên bản mới; trên dữ liệu đầy đủ có 96.096 khách và thêm 242 phiên bản (khoảng 0,25\%), nên chi phí lưu lịch sử thấp. Chiều người bán \code{dim\_\allowbreak{}seller} là chiều dùng chung (conformed) nối quy trình bán hàng với quy trình thu hút seller, dùng SCD loại 1; 462 seller chỉ có trong CRM vẫn có dòng với địa chỉ chưa biết. Chiều sản phẩm gom 73 danh mục vào 13 nhóm ngành hàng và phân loại cỡ hàng theo cân nặng. Các thuộc tính nhỏ của deal được gom vào một chiều rác (junk dimension) gồm 503 tổ hợp, và nhân viên sales được tách thành hai vai trò SDR và SR trong cùng một bảng. Toàn bộ 39 phép kiểm thử dimension đều đạt: mỗi chiều có đúng một dòng Unknown, mỗi khách có đúng một phiên bản hiện hành và các khoảng hiệu lực không chồng lấn.

\begin{table}[htbp]
\centering\small
\caption{Các bảng chiều: grain, số dòng (gồm dòng Unknown $-1$) và loại SCD}
\label{tab:dims}
\begin{tabular}{llrc}
\toprule
\textbf{Chiều} & \textbf{Grain} & \textbf{Số dòng} & \textbf{SCD} \\
\midrule
dim\_date & 1 ngày (2016-01-01 → 2020-12-31) & 1.828 & -- \\
dim\_geography & 1 mã bưu chính (zip) & 19.178 & 1 \\
dim\_customer & 1 phiên bản địa chỉ của 1 khách thật & 96.339 & 2 \\
dim\_seller & 1 seller (chiều dùng chung) & 3.558 & 1 \\
dim\_product & 1 sản phẩm & 32.952 & 1 \\
dim\_order\_status & 1 trạng thái đơn & 9 & 1 \\
dim\_payment\_type & 1 hình thức thanh toán & 6 & 1 \\
dim\_lead\_origin & 1 giá trị origin của MQL & 11 & 1 \\
dim\_lead\_profile & 1 tổ hợp thuộc tính deal (junk dimension) & 503 & -- \\
dim\_sales\_rep & 1 nhân viên và vai trò (SDR/SR) & 55 & -- \\
\bottomrule
\end{tabular}
\end{table}


\subsection{Hệ thống chiều khái niệm}\label{sec:he-thong-chieu}
Mỗi chiều có một hoặc hai phân cấp; các bảng dưới đây liệt kê từng cấp thành một cột, vài giá trị ví dụ thật ở các dòng giữa và số giá trị phân biệt của từng cấp ở dòng cuối (đếm từ kho, không tính dòng Unknown). Chiều thời gian đi từ năm xuống ngày (5 năm, 20 quý, 60 tháng, 1.827 ngày); chiều địa lý từ vùng xuống mã bưu chính (5 vùng, 27 bang, 6.190 cặp bang--thành phố, 19.177 mã); chiều sản phẩm từ 13 nhóm ngành xuống 74 danh mục và 32.951 sản phẩm; chiều seller dùng cùng phân cấp địa lý (3.557 seller) và thêm phân cấp kênh thu hút; chiều khách hàng thêm cấp khách thật (96.096) và phiên bản địa chỉ (96.338 dòng) vì dùng SCD loại 2.
\begin{table}[htbp]\centering\scriptsize\caption{Chiều thời gian: Năm → Quý → Tháng → Ngày}\label{tab:hier-dim-date}
\begin{tabular}{llll}\toprule
\textbf{Năm} & \textbf{Quý} & \textbf{Tháng} & \textbf{Ngày}\\\midrule
2016 & 2016-Q4 & 2016-12 & 2016-12-25\\
2017 & 2017-Q4 & 2017-11 & 2017-11-24\\
2018 & 2018-Q1 & 2018-03 & 2018-03-15\\
\midrule
\textbf{5} & \textbf{20} & \textbf{60} & \textbf{1.827}\\
\multicolumn{4}{l}{\footnotesize\emph{Dòng cuối: số giá trị phân biệt của từng cấp (không tính dòng Unknown).}}\\
\bottomrule\end{tabular}\end{table}

\begin{table}[htbp]\centering\scriptsize\caption{Chiều địa lý: Vùng → Bang → Thành phố → Mã bưu chính}\label{tab:hier-dim-geography}
\begin{tabular}{llll}\toprule
\textbf{Vùng} & \textbf{Bang} & \textbf{Thành phố (kèm bang)} & \textbf{Mã bưu chính}\\\midrule
Nordeste & BA & salvador & 40010\\
Sul & RS & porto alegre & 90001\\
Sudeste & SP & sao paulo & 01001\\
\midrule
\textbf{5} & \textbf{27} & \textbf{6.190} & \textbf{19.177}\\
\multicolumn{4}{l}{\footnotesize\emph{Dòng cuối: số giá trị phân biệt của từng cấp (không tính dòng Unknown).}}\\
\bottomrule\end{tabular}\end{table}

\begin{table}[htbp]\centering\scriptsize\caption{Chiều sản phẩm: Nhóm ngành → Danh mục → Sản phẩm}\label{tab:hier-dim-product}
\begin{tabular}{lll}\toprule
\textbf{Nhóm ngành hàng} & \textbf{Danh mục (tiếng Anh)} & \textbf{Sản phẩm (product\_id)}\\\midrule
Electronics \& Computers & computers\_accessories & 0021a87d\\
Health \& Beauty & health\_beauty & 00210e41\\
Home \& Furniture & bed\_bath\_table & 0009406f\\
\midrule
\textbf{13} & \textbf{74} & \textbf{32.951}\\
\multicolumn{3}{l}{\footnotesize\emph{Dòng cuối: số giá trị phân biệt của từng cấp (không tính dòng Unknown).}}\\
\bottomrule\end{tabular}\end{table}

\begin{table}[htbp]\centering\scriptsize\caption{Chiều seller: Vùng → Bang → Thành phố → Seller}\label{tab:hier-dim-seller}
\begin{tabular}{llll}\toprule
\textbf{Vùng} & \textbf{Bang} & \textbf{Thành phố (kèm bang)} & \textbf{Seller}\\\midrule
Sudeste & MG & belo horizonte & 00d8b143\\
Sul & PR & curitiba & 01266d4c\\
Sudeste & SP & sao paulo & 00ab3eff\\
\midrule
\textbf{6} & \textbf{23} & \textbf{636} & \textbf{3.557}\\
\multicolumn{4}{l}{\footnotesize\emph{Dòng cuối: số giá trị phân biệt của từng cấp (không tính dòng Unknown).}}\\
\bottomrule\end{tabular}\end{table}

\begin{table}[htbp]\centering\scriptsize\caption{Chiều khách hàng (SCD loại 2): Vùng → Bang → Thành phố → Mã bưu chính → Khách → Phiên bản}\label{tab:hier-dim-customer}
\begin{tabular}{llllll}\toprule
\textbf{Vùng} & \textbf{Bang} & \textbf{Thành phố (kèm bang)} & \textbf{Mã bưu chính} & \textbf{Khách thật} & \textbf{Phiên bản (dòng)}\\\midrule
Sudeste & MG & belo horizonte & 31555 & 000c8bdb & 1\\
Sudeste & RJ & rio de janeiro & 21380 & 00293787 & 1\\
Sudeste & SP & sao paulo & 03071 & 001a2bf0 & 1\\
\midrule
\textbf{5} & \textbf{27} & \textbf{4.310} & \textbf{14.994} & \textbf{96.096} & \textbf{96.338}\\
\multicolumn{6}{l}{\footnotesize\emph{Dòng cuối: số giá trị phân biệt của từng cấp (không tính dòng Unknown).}}\\
\bottomrule\end{tabular}\end{table}

\begin{table}[htbp]\centering\scriptsize\caption{Chiều kênh thu hút: Nhóm kênh → Origin}\label{tab:hier-dim-lead-origin}
\begin{tabular}{ll}\toprule
\textbf{Nhóm kênh} & \textbf{Origin}\\\midrule
Email & email\\
Organic & direct\_traffic\\
Other & other\\
\midrule
\textbf{6} & \textbf{10}\\
\multicolumn{2}{l}{\footnotesize\emph{Dòng cuối: số giá trị phân biệt của từng cấp (không tính dòng Unknown).}}\\
\bottomrule\end{tabular}\end{table}

\begin{table}[htbp]\centering\scriptsize\caption{Chiều trạng thái đơn: Nhóm trạng thái → Trạng thái}\label{tab:hier-dim-order-status}
\begin{tabular}{ll}\toprule
\textbf{Nhóm trạng thái} & \textbf{Trạng thái đơn}\\\midrule
Completed & delivered\\
Failed & canceled\\
In progress & approved\\
\midrule
\textbf{3} & \textbf{8}\\
\multicolumn{2}{l}{\footnotesize\emph{Dòng cuối: số giá trị phân biệt của từng cấp (không tính dòng Unknown).}}\\
\bottomrule\end{tabular}\end{table}

\begin{table}[htbp]\centering\scriptsize\caption{Chiều hình thức thanh toán: Nhóm → Hình thức}\label{tab:hier-dim-payment-type}
\begin{tabular}{ll}\toprule
\textbf{Nhóm thanh toán} & \textbf{Hình thức}\\\midrule
Bank slip & boleto\\
Card & credit\_card\\
Unknown & not\_defined\\
Voucher & voucher\\
\midrule
\textbf{4} & \textbf{5}\\
\multicolumn{2}{l}{\footnotesize\emph{Dòng cuối: số giá trị phân biệt của từng cấp (không tính dòng Unknown).}}\\
\bottomrule\end{tabular}\end{table}

\begin{table}[htbp]\centering\scriptsize\caption{Chiều nhân viên sales: Vai trò → Nhân viên}\label{tab:hier-dim-sales-rep}
\begin{tabular}{ll}\toprule
\textbf{Vai trò} & \textbf{Nhân viên}\\\midrule
SDR & 068066e2\\
SR & 060c0a26\\
\midrule
\textbf{2} & \textbf{54}\\
\multicolumn{2}{l}{\footnotesize\emph{Dòng cuối: số giá trị phân biệt của từng cấp (không tính dòng Unknown).}}\\
\bottomrule\end{tabular}\end{table}


\subsection{Kiểm tra tính đúng của các phân cấp}\label{sec:kiem-tra-phan-cap}
Một phân cấp đúng khi mỗi giá trị ở cấp dưới thuộc đúng một giá trị ở cấp trên. Bảng \ref{tab:hier-checks} ghi kết quả đếm trên kho cho các phân cấp dễ sai nhất. \textbf{Tháng phải đi kèm năm}: chỉ có 12 số tháng nhưng 60 cặp (năm, tháng), nên cấp tháng dùng \code{year\_month} chứ không dùng số tháng một mình. \textbf{Thành phố phải đi kèm bang}: có 5.814 tên thành phố nhưng 6.190 cặp (bang, thành phố), và 303 tên thành phố xuất hiện ở từ hai bang trở lên (nhiều nhất \code{bom jesus} và \code{sao domingos}, mỗi tên ở 6 bang); vì vậy cấp thành phố luôn gắn với bang và không gộp theo tên. Các phân cấp còn lại đều là quan hệ một nhiều: mỗi mã bưu chính thuộc một thành phố và một bang, mỗi bang thuộc một vùng, mỗi danh mục thuộc một nhóm ngành và mỗi origin thuộc một nhóm kênh.
\begin{table}[htbp]\centering\scriptsize\caption{Kiểm tra tính đúng của các phân cấp (đếm từ kho dữ liệu)}\label{tab:hier-checks}
\begin{tabular}{p{3.6cm}p{7.2cm}p{4.6cm}}\toprule
\textbf{Phân cấp} & \textbf{Kết quả đếm} & \textbf{Kết luận}\\\midrule
Tháng phải đi kèm năm & chỉ có 12 số tháng nhưng 60 cặp (năm, tháng): tên tháng một mình không xác định một nút của phân cấp & đúng: dùng year\_month ('YYYY-MM') và year\_quarter làm cấp, không dùng month một mình\\
Thành phố phải đi kèm bang & 5.814 tên thành phố nhưng 6.190 cặp (bang, thành phố); 303 tên thành phố xuất hiện ở từ 2 bang trở lên (nhiều nhất: bom jesus (6 bang), sao domingos (6 bang), boa esperanca (5 bang)) & đúng: cấp thành phố luôn đi kèm bang; không gộp theo tên thành phố\\
Mã bưu chính xác định một thành phố và một bang & 0 mã có nhiều hơn một thành phố, 0 mã có nhiều hơn một bang & đúng: dim\_geography có một dòng cho mỗi mã bưu chính\\
Bang thuộc đúng một vùng & 0 bang thuộc nhiều hơn một vùng (27 bang, 5 vùng) & đúng: Vùng → Bang là quan hệ một nhiều\\
Sản phẩm → danh mục → nhóm ngành & 0 danh mục thuộc nhiều hơn một nhóm ngành (74 danh mục, 13 nhóm ngành) & đúng: mỗi danh mục một nhóm; sản phẩm chưa có danh mục vào nhóm unknown\\
Origin → nhóm kênh & 0 origin thuộc nhiều hơn một nhóm kênh (10 origin, 6 nhóm kênh) & đúng: mỗi origin một nhóm kênh\\
\bottomrule\end{tabular}\end{table}


\subsection{Phân nhóm các dimension}\label{sec:phan-nhom-dim}
Mười bảng chiều được chia theo vai trò trong mô hình (Bảng \ref{tab:dim-groups}): \emph{chiều dùng chung} (conformed) là chiều được từ hai bảng fact trở lên dùng chung, gồm \code{dim\_date} (6 fact), \code{dim\_customer} (4), \code{dim\_seller} (3), \code{dim\_geography} và \code{dim\_order\_status} (2); \emph{chiều đóng vai} (role-playing) là \code{dim\_date} (11 view cho các mốc mua, duyệt, bàn giao, giao, dự kiến, liên hệ, chốt) và \code{dim\_sales\_rep} (hai vai SDR và SR); \emph{chiều SCD loại 2} là \code{dim\_customer}; \emph{chiều rác} là \code{dim\_lead\_profile} (503 tổ hợp của 7 thuộc tính của deal); \emph{danh mục nhỏ} là các bảng tra cứu 6 đến 11 dòng; và \code{dim\_product} là chiều của một quy trình (chỉ \code{fact\_order\_item} dùng).
\begin{table}[htbp]\centering\scriptsize\caption{Phân nhóm 10 bảng chiều (một chiều có thể thuộc nhiều nhóm)}\label{tab:dim-groups}
\begin{tabular}{p{3.6cm}p{3.0cm}rrp{5.0cm}}\toprule
\textbf{Nhóm} & \textbf{Chiều} & \textbf{Số dòng} & \textbf{Số fact} & \textbf{Ghi chú}\\\midrule
Chiều dùng chung (conformed) & \code{dim\_date} & 1.828 & 6 & 11 view đóng vai; dùng bởi mọi fact\\
Chiều dùng chung (conformed) & \code{dim\_geography} & 19.178 & 2 & zip → thành phố → bang → vùng; nối khách, seller\\
Chiều dùng chung (conformed) & \code{dim\_seller} & 3.558 & 3 & nối quy trình bán hàng với thu hút seller (SCD1)\\
Chiều dùng chung (conformed) & \code{dim\_customer} & 96.339 & 4 & SCD loại 2: giữ lịch sử mã bưu chính\\
Chiều dùng chung (conformed) & \code{dim\_order\_status} & 9 & 2 & dùng bởi fact đơn hàng và fact xử lý đơn\\
Chiều đóng vai (role-playing) & \code{dim\_date} & 1.828 & 6 & mua, duyệt, bàn giao, giao, dự kiến, liên hệ, chốt, đơn đầu tiên...\\
Chiều đóng vai (role-playing) & \code{dim\_sales\_rep} & 55 & 1 & một bảng cho hai vai trò: SDR và SR\\
Chiều SCD loại 2 & \code{dim\_customer} & 96.339 & 4 & đóng phiên bản cũ, mở phiên bản mới khi đổi mã bưu chính\\
Chiều rác (junk) & \code{dim\_lead\_profile} & 503 & 1 & gom 7 thuộc tính nhỏ của deal thành một chiều\\
Danh mục nhỏ (lookup) & \code{dim\_order\_status} & 9 & 2 & trạng thái đơn → nhóm trạng thái\\
Danh mục nhỏ (lookup) & \code{dim\_payment\_type} & 6 & 1 & hình thức thanh toán → nhóm\\
Danh mục nhỏ (lookup) & \code{dim\_lead\_origin} & 11 & 1 & origin → nhóm kênh\\
Chiều của một quy trình (một fact) & \code{dim\_product} & 32.952 & 1 & sản phẩm → danh mục → nhóm ngành; chỉ fact\_order\_item dùng\\
\bottomrule\end{tabular}\end{table}

\todo{Nhóm vẽ lại hệ thống chiều khái niệm (cây phân cấp từng chiều) bằng draw.io từ các bảng trên; bus matrix ở docs/DW\_DESIGN.md mục 4.}
